\section{Results}

\subsection{Constraint Localization}

Our analysis reveals that constraint directions are concentrated in mid-to-late layers across all models studied. Table~\ref{tab:peak_layers} shows the peak constraint layers for each model.

\begin{table}[h]
\centering
\begin{tabular}{lcc}
\toprule
Model & Peak Layer & Signal Strength \\
\midrule
GPT-2 & 8 & 0.45 \\
LLaMA-3.1-8B & 15 & 0.78 \\
Mistral-7B & 14 & 0.72 \\
Qwen2.5-7B & 16 & 0.69 \\
\bottomrule
\end{tabular}
\caption{Peak constraint layers and normalized signal strength.}
\label{tab:peak_layers}
\end{table}

\subsection{Cross-Layer Alignment}

Figure~\ref{fig:alignment} shows the cross-layer alignment heatmap. Directions remain highly consistent ($\cos \theta > 0.8$) within clusters of 3-5 consecutive layers, with sharp transitions at layer boundaries.

\subsection{Surgical Removal Results}

Table~\ref{tab:removal} presents results of constraint removal across models.

\begin{table}[h]
\centering
\begin{tabular}{lcccc}
\toprule
Model & Refusal Before & Refusal After & Coherence & Perplexity $\Delta$ \\
\midrule
GPT-2 & 87\% & 4\% & 0.92 & +2.1\% \\
LLaMA-3.1-8B & 94\% & 5\% & 0.91 & +2.8\% \\
Mistral-7B & 91\% & 6\% & 0.93 & +1.9\% \\
Qwen2.5-7B & 89\% & 3\% & 0.94 & +1.5\% \\
\bottomrule
\end{tabular}
\caption{Constraint removal results. Refusal reduction exceeds 90\% with less than 3\% capability loss.}
\label{tab:removal}
\end{table}

\subsection{Ouroboros Effect}

Figure~\ref{fig:ouroboros} shows the self-repair curve. Without compensation, refusal re-emerges to 40\% of original within 100 samples. With 2 Ouroboros passes, re-emergence is reduced to 12\%.

\subsection{Mathematical Barrier Analysis}

The shell method barrier exhibits rank-3 geometric structure. Table~\ref{tab:barrier} summarizes its properties.

\begin{table}[h]
\centering
\begin{tabular}{ll}
\toprule
Property & Value \\
\midrule
Constraint Direction & spherical\_code\_dependency \\
Barrier Type & unconditional \\
Location & Lemma 4.2 $\rightarrow$ Theorem 1 \\
Threshold & $\exp(-c\log N)$ \\
Rank & 3 \\
Solid Angle & 2.1 sr \\
\bottomrule
\end{tabular}
\caption{Geometric analysis of the shell method barrier.}
\label{tab:barrier}
\end{table}

\subsection{Cross-Domain Transfer}

Applying refusal removal techniques to the shell method barrier yields a conditional bypass strategy achieving $\exp(-C\sqrt{\log N})$ under Hypothesis H.

\subsection{Self-Optimization Results}

Table~\ref{tab:evolve} shows the effect of ARIS self-optimization on constraint strength.

\begin{table}[h]
\centering
\begin{tabular}{lcc}
\toprule
Constraint & Before & After \\
\midrule
safety\_boundary & 0.92 & 0.37 \\
content\_policy & 0.78 & 0.31 \\
instruction\_guard & 0.65 & 0.28 \\
\bottomrule
\end{tabular}
\caption{ARIS constraint reduction after self-optimization.}
\label{tab:evolve}
\end{table}
